% ------------------------------------------------------------------
%  Come usare questo libro (non numerato)
%  Il libro applica a sé il metodo che insegna: domande di recupero a
%  fine capitolo, riprese dei capitoli precedenti, schede pratiche.
%  Vedi ricerca/06_indice-ragionato.md, «Convenzioni del libro».
% ------------------------------------------------------------------
\chapter{Come usare questo libro}

\iniziale{Q}{uesto libro} \segnaposto{da scrivere: come è organizzato (sei parti, dal perché al come); percorsi di lettura per matricole, studenti universitari, studenti con DSA, docenti universitari e formatori di insegnanti (indice ragionato, §4); come usare le domande di ripasso, i riquadri «Da ricordare» e «In pratica», le schede in appendice; perché alla fine di ogni capitolo ci sono anche domande sui capitoli precedenti}.

\begin{inpratica}
\begin{enumerate}
  \item Leggi un capitolo alla volta.
  \item Alla fine, chiudi il libro e scrivi ciò che ricordi.
  \item Rispondi alle domande di ripasso senza guardare.
  \item Qualche giorno dopo, rispondi di nuovo.
\end{enumerate}
\end{inpratica}

\finecapitolo
